% measure-dissolution.tex — paper/measure-dissolution-full.md の LaTeX 化 (v18.x 出版作業)
% 対象誌: JHEP / PRD (letter 級)。数値の一次ソースは results/ — 本文数値は .md 原稿から転記
% TODO(著者確認): 著者名・所属は仮置き。投稿前に確定すること。
\documentclass[aps,prd,preprint,nofootinbib]{revtex4-2}
\usepackage[utf8]{inputenc}
\usepackage[T1]{fontenc}
\usepackage{amsmath,amssymb}
\usepackage{booktabs}
\usepackage{hyperref}

\begin{document}

\title{The Measure Problem That Dissolved:\\
A Pre-Registered Contest of Data-Blind Priors over Magnetized-Torus Yukawa Models}

\author{Yasuo Higano}
\affiliation{Independent researcher; Quantum Relational Network program}

\date{\today}

\begin{abstract}
If the discrete data of a compactification --- fluxes, Wilson lines, generation
pairing, complex structure --- are irreducible, the residual scientific
question is the \emph{measure}: which data-blind prior should weigh the model
space? We report a complete, pre-registered arc on this question, run over the
magnetized-torus Yukawa program of the companion papers. A hierarchy-tilted
prior, catastrophically rejected on the rectangular geometry, becomes the
program's first surviving selection principle on the CP-capable shear family
($+2.3$ nats); its anatomy (survival window $\beta\in[0.25,1.7]$, collapse of
hard selection at $-8386$ nats) and identity (bin-free and binned
density-flattening priors reproduce its peak within 0.05 nats) reduce it to a
\emph{measure correction}: the configuration population over-represents shallow
mass ratios. We then registered six candidate measures with three success
criteria fixed in advance --- data-blind construction, convention invariance,
and non-degradation of never-fitted holdouts --- together with a continuum
marginalization of the Higgs width (removing the last prior-grid convention)
and a moment-based analytic form (which succeeds exactly on the geometry family
whose configuration density is moment-compatible, and fails on the heavy-tailed
diagonal family). The judgment kills the entire correction class: the only
candidate to pass the evidence and stability criteria degrades the
unitarity-triangle holdout and is rejected by the third criterion, whose first
live use this is. The cause is a lever the corrections had been compensating
for: scanning the \emph{imaginary} part of the complex structure finds an
interior evidence valley at $\tau=1/12+i/2$, where the ten-observable evidence
reaches the program's best value, the two chronic tensions (a two-fold Cabibbo
excess and a low Jarlskog invariant) resolve simultaneously at MAP (factors
1.01 and 1.24) and at the posterior-band level, all twelve observables sit
within a factor of five with six at the ten-percent level --- \textbf{and every
measure correction is now strictly harmful}. The measure problem, posed
properly, dissolved into geometry: the uniform counting measure is correct once
$\tau$ is right, and what remains is a thermodynamic prior's stubborn,
sub-threshold positive signal ($+0.3$ to $+0.6$ nats, sign-stable across four
generations of windows) --- the sharpest surviving hint of a dynamical
selection principle.
\end{abstract}

\maketitle

\section{Introduction}

The companion papers~\cite{Companion,CPpaper} closed with a deferral: the
discrete data of the compactification ``await a common selection principle.''
The first attempt (nine observables, rectangular geometry) rejected or left
undecided every candidate --- minimum description length, self-stability,
hierarchy maximization (at $-494$ nats), thermodynamic condensation ($+0.32$,
the lone positive sign). This paper reports what happened when the question was
asked again, properly: on the CP-capable geometry family of~\cite{CPpaper},
with the Jarlskog invariant in the likelihood, with candidates and success
criteria registered before the runs, and with every instrument convention
(prior grids, histogram bins) either eliminated or promoted to a tested
hypothesis.

The arc has a beginning (a principle survives for the first time), a middle
(its identity is found, its class is registered and judged), and an end we did
not anticipate: the measure problem dissolves. The discipline that carried the
companion papers --- pre-registered branch verdicts, regression anchors as
gates, negative results reported with the same machinery as positive ones ---
is what made the dissolution \emph{visible}: had we tuned candidates
window-by-window, the compensation structure described in \S\ref{sec:dissolve}
would have read as success.

\section{A principle survives, and is dissected}

On the 21-geometry shear family (both tori sheared by $s_1,s_2\in\{0..5\}$ at
lattice size $N=36$), with the ten-observable evidence of~\cite{CPpaper}, the
four principles of the first contest were re-run as two-level data-blind
priors. Three verdicts repeated (MDL now \emph{inherits the CP penalty} through
its preference for the rectangular geometry, $-2.16$; self-stability rejected
again; thermodynamic again positive, $+0.41$). One reversed spectacularly: the
hierarchy-tilted prior $e^{\beta\cdot\mathrm{depth}}$, worst of the first
contest, \textbf{survives at $+2.31$ nats} --- the program's first surviving
selection principle. Its per-geometry anatomy explains the reversal: on
symmetric (equal-shear) geometries its deep tail overshoots the observed
hierarchy exactly as on the rectangle ($-514.9$, reproducing the first
contest's $-514.4$), while on asymmetric geometries the tilt concentrates prior
mass where the likelihood lives.

Dissection followed. The survival curve over $\beta$ peaks at $\beta=1.0$
($\Delta=+3.34$), collapses below the survival threshold at $\beta^*=2.0$, and
the analytic $\beta\to\infty$ endpoint --- hard selection of each geometry's
deepest configuration --- is a catastrophe at $-8386$ nats. Asymmetric tails
are two orders of magnitude shallower than symmetric ones, but both overshoot
the observation. ``The universe does not select the deepest'' survives the
change of family; \textbf{what survives is a tilt of order one.}

\section{The survivor is a measure correction}

Since $e^{\beta\cdot\mathrm{depth}}=\prod_i r_i^{-\beta}$, the surviving tilt
reads as a conversion toward a scale-invariant (log-uniform,
Jeffreys-type~\cite{Jeffreys}) measure on mass ratios. The direct test builds
data-blind histogram-flattening priors from the configuration population itself
(1D in depth; 2D in the two log-ratios; three bin widths): flattening
reproduces the peak within \textbf{0.05 nats}, all variants landing in a
$+2.9\ldots+3.3$ band, while the naive mechanism is refuted --- the
configuration density is \emph{not} exponential in depth (measured slope
$-0.19$, not $-1$). The operative content is a \textbf{measure correction}: the
configuration measure of the shear family over-represents shallow mass ratios,
and removing this over-representation --- by tilt, by flattening, in one or two
dimensions --- is worth $+3$ nats, with no observation consulted.

\section{Registration and instruments}

At this point the program registered the question properly. Six candidate
measures (uniform baseline; Jeffreys/Fisher; MDL; thermodynamic; depth tilt,
$\beta$ marginalized; histogram flattening) were entered in a ledger with three
success criteria fixed before any judgment run:
\begin{itemize}
\item \textbf{S1 (data-blind)}: constructible without observations, \emph{with
      every convention --- bins, grids --- counted as part of the
      construction};
\item \textbf{S2 (convention invariance)}: candidate rankings survive changes
      of bin width, prior grid, lattice size, and geometry window;
\item \textbf{S3 (holdout non-degradation)}: a surviving measure must not
      worsen never-fitted observables, registered in the prediction ledger
      before the run (the unitarity-triangle angles $\beta$ and $\gamma$ here).
\end{itemize}

Two instruments were built to make the criteria testable. First, the
Higgs-width prior grid --- the one convention shown earlier to flip sub-nat
rankings --- was replaced by a continuum marginalization (trapezoid over
$\sigma\in[1.2,8.0]$; the convergence gate failed at the first grid density and
was passed by doubling it, with the failure recorded). The marginalization
confirmed that two earlier conclusions were not grid artifacts, and
re-confirmed the program's most robust continuous structure: the evidence
concentrates at physical Higgs width $\sigma/\sqrt{\mathrm{area}}=1/18$ across
every geometry, lattice, and grid tested. Second, a \emph{bin-free} analytic
form for the flattening measure was constructed from sample moments (a
quadratic tilt $w=\tfrac12(x-\mu)^{\!\top}\Sigma^{-1}(x-\mu)$, zero fitting
freedom). It succeeds exactly where moments exist: on the asymmetric family the
depth marginal is near-Gaussian (quantile-quantile RMS 0.25) and the quadratic
tilt reproduces binned flattening to 0.10 nats --- S1 fully satisfied, bins
eliminated. On the diagonal family the sample variance is dominated by the
anomalous deep tail (QQ $\approx 52$) and the quadratic tilt destroys itself
($-350$ nats) --- the same tail that killed hard selection, now killing moment
estimation. \emph{Whether the measure correction has an analytic form is itself
a property of the geometry family.}

\section{The judgment}

All candidates were then judged simultaneously, on three windows (asymmetric,
diagonal, and the best-evidence rectangular-lattice geometry then known), under
the continuum $\sigma$ marginalization, against the pre-registered criteria.
\textbf{Survivors: none.}
\begin{itemize}
\item MDL: rejected a third time ($-1.64$, sign-stable).
\item Depth: rejected --- the diagonal windows' tails collapse it (S2
      sign-unstable: $+2.84$ vs $-494$).
\item Flattening (1D and 2D): rejected --- \textbf{the best-geometry window
      reverses its advantage} ($-2.15$ and $-0.38$), breaking S2.
\item The bin-free quadratic tilt: passes evidence ($+1.93$) and S2 on its
      qualified window, and is \textbf{rejected by S3} --- it degrades the
      unitarity-triangle holdout from 1.96 to 4.53 rad. This is S3's first live
      kill, and the sharpest lesson of the arc: \emph{a measure that raises the
      evidence and a measure that improves predictions are different objects.}
\item Thermodynamic: undecided at $+0.49$ --- but sign-positive in \emph{every}
      window, as it has been in every era ($+0.32$, $+0.41$, $+0.49$).
\end{itemize}

\section{The dissolution: geometry was the missing variable}
\label{sec:dissolve}

The window that broke the flattening class was the clue. The aspect ratio of
the torus --- the \emph{imaginary} part of the complex structure, fixed at 1 in
every previous version --- turned out to be a live lever: scanning it found
monotone improvement toward wider tori and then an interior valley bottom at
$\boldsymbol{\tau=1/12+i/2}$ (lattice $36\times 18$, diagonal shear $(3,3)$,
physical Higgs width $1/18$). At this point the ten-observable evidence reaches
the program's best value ($-18.43$, beating the rectangular nine-observable
champion by 1.4 nats on the harder scale); the two chronic tensions resolve
\emph{simultaneously} --- the Cabibbo element, once overpredicted by 2.5, sits
at factor 1.01, and the Jarlskog invariant at 1.24, with both 68\% posterior
bands covering the measurements for the first time; the full scorecard holds
all twelve observables within a factor of five, six of them at the ten-percent
level.

And at this point, \textbf{every measure correction is strictly harmful} (1D
flattening $-3.73$; 2D $-1.23$; the quadratic tilt disqualified and
self-destructive; the depth tilt $-1004$). Re-examined across the
$\tau_{\rm im}$ scan, the pattern is monotone: the better the geometry, the
less the correction helps, until it hurts. The $+3$ nats that the correction
class earned in \S3 was \emph{compensation for a mis-set modulus} --- the
configuration measure of the $\tau_{\rm im}=1$ section over-represents shallow
ratios precisely because that section's geometry cannot reach the observed
hierarchy pattern efficiently. Fix $\tau$, and the uniform counting measure is
not merely adequate; it is preferred.

We record the arc's conclusion in the form the ledger demands: \textbf{the
measure problem, as posed, dissolved into geometry.} The surviving question is
no longer ``which prior on configurations'' but ``which prior on $\tau$'' ---
one level up, where the program's evidence now concentrates on a specific
rational point ($\tau_{\rm re}=1/12$ from one scan, $\tau_{\rm im}=1/2$ from
another) whose arithmetic ($n_x=2n_y$; exact three-fold degeneracy at every
aspect ratio tested, to $10^{-13}$) belongs to the same open lattice number
theory as the companion's index conditions.

\section{What survives: the thermodynamic whisper}

One signal refused to die at any stage: the thermodynamic prior (weight
$\propto e^{\gamma\ln\|Y\|_F}$, a minimal proxy for condensation energy) is
positive in every window of every era --- $+0.32$ (rectangular, nine
observables), $+0.41$ (shear family, ten), $+0.44/+0.40/+0.63$ (the three
judgment windows), $+0.51$ (the $\tau$ valley bottom) --- and never reaches the
$+1$ survival line. Four generations of windows, one sign, no verdict. It is
the only candidate whose content is \emph{dynamical} rather than statistical
(it rewards strong coupling, not deep hierarchy), and its registered upgrade
--- replacing the Frobenius proxy by a one-loop vacuum energy with the
Kaluza--Klein tower --- is now the sharpest open item on the
selection-principle ledger.

\section{Limitations}

\begin{itemize}
\item All judgments are conditioned on the likelihood conventions of the
      companion program (lognormal, $\sigma=\ln 2$; ten observables) and on the
      shear-family model space; ``dissolved'' means \emph{within this program's
      model space and observables}, not as a theorem.
\item The $\tau$ valley was located by two one-dimensional scans
      ($\tau_{\rm re}$, then $\tau_{\rm im}$); the two-dimensional landscape
      and the arithmetic of its rational points are unexplored.
\item The residual weak spots at the valley bottom are $m_c/m_t$ (factor 2.04
      --- the companion's historic weak spot, untouched by the $\tau$ lever), a
      mild high bias of the $|V_{cb}|/|V_{ts}|$ posterior bands, and a global
      sign: all CP-odd angles come out with magnitudes near measurement but
      negative sign, suggesting an orientation ($\tau\to-\tau$) test that the
      $|J|$-based likelihood cannot see. [The orientation test was subsequently
      performed: the anti-unitary branch pair is exact to machine precision,
      the measured $\mathrm{sign}(J)>0$ fixes the branch, and the never-fitted
      $\gamma=+66.8^\circ$ then falls within the experimental uncertainty of
      $65.9^\circ\pm 3.5^\circ$.]
\item The thermodynamic candidate's proxy is not yet a one-loop energy; its
      persistence could still be an artifact of the proxy's correlation with
      overlap magnitudes.
\end{itemize}

\section{Reproducibility}

Deterministic and dependency-free as in the companions (Rust std-only; built-in
PASS/FAIL gates; machine-validated claim ledger; 469 PASS / 0 FAIL at the
period-17 integration preceding this arc). Practices specific to this arc: (i)
a \emph{measure ledger} (measures.yml) registering candidates, conventions, and
success criteria before judgment; (ii) \emph{holdout-based rejection} (S3)
wired into the same run that freezes the holdout baseline before scoring ---
the ordering is machine-enforced in the program's stdout; (iii) a
\emph{mode-table disk cache} with explicitly non-primary status, which cut the
marginal cost of a full judgment (fifteen prior variants $\times$ three windows
$\times$ 69 Higgs widths $\times$ 1.7M-configuration sums) to about one minute;
(iv) convergence gates on numerical integration, with a recorded instance of
the gate failing and being passed by refinement rather than relaxation.

\begin{thebibliography}{9}
\bibitem{Companion} This program: ``Yukawa Hierarchies from Magnetized Tori without Order-One Coefficients,'' paper/geometric-yukawa-full.md (LaTeX version: paper/tex/geometric-yukawa.tex).
\bibitem{CPpaper} This program: ``CP Violation Requires Complex Structure,'' paper/cp-complex-structure-full.md (LaTeX version: paper/tex/cp-complex-structure.tex).
\bibitem{Jeffreys} H.~Jeffreys, ``An invariant form for the prior probability in estimation problems,'' Proc.\ R.\ Soc.\ Lond.\ A \textbf{186} (1946) 453--461.
\bibitem{Trotta} R.~Trotta, ``Bayes in the sky: Bayesian inference and model selection in cosmology,'' Contemp.\ Phys.\ \textbf{49} (2008) 71--104.
\bibitem{Repo} This program's repository: Quantum Relational Network (\texttt{claims.yml}, \texttt{measures.yml}, \texttt{predictions.yml}, \texttt{results/}) --- primary source for all numbers.
\end{thebibliography}

\end{document}
